% ============================================================================
% DRAFT-READY SECTION: The Resting Tension of the Planck Field
% For insertion into universal_condensate_draft_v02.tex
% Author: Math-Agent
% Date: \today
% ============================================================================
%
% USAGE: Insert this section after Section 3.4 (The Density Hierarchy)
% or as a new subsection within Section 3 (From Gross-Pitaevskii to FLRW).
%
% Prerequisites: Requires the macros \rhostiff, \rhocrit, \rhoPl, \lPl, \mPl, \RH
% already defined in the main document preamble.
%
% This section is approximately 2-3 pages when compiled.
% ============================================================================

\subsection{The Resting Tension of the Planck Field}
\label{sec:resting_tension}

The stiffness-matching condition (Eq.~\ref{eq:stiffness_match}) defines a fundamental density scale:
\begin{equation}
\label{eq:rho_stiff_tension}
\rhostiff = \frac{c^2}{8\pi G} = 5.37 \times 10^{25} \text{ kg/m}^3.
\end{equation}

This density is $\sim 10^{52}$ times larger than the cosmological critical density $\rhocrit \approx 9.47 \times 10^{-27}$ kg/m$^3$. Rather than viewing this ratio as a discrepancy to be resolved, we recognize that \textbf{$\rhostiff$ and $\rhocrit$ are fundamentally different physical quantities}, and their comparison reveals the physical mechanism underlying the weakness of gravity.

\subsubsection{Constitutive vs. Content Properties}
\label{sec:const_vs_content}

In continuum mechanics, \emph{constitutive properties} describe intrinsic material characteristics---how a medium responds to stress or deformation---independent of external contents. Examples include Young's modulus $E$, bulk modulus $K$, and tensile strength. In contrast, \emph{content properties} describe the amount or distribution of matter residing within a medium.

\begin{table}[H]
\centering
\caption{Constitutive vs. Content Properties: Physical Analogies}
\label{tab:analogies}
\begin{tabular}{@{}lll@{}}
\toprule
\textbf{System} & \textbf{Constitutive Property} & \textbf{Content Property} \\
\midrule
Guitar string & String tension $T$ & Mass of objects hanging from string \\
Steel table & Elastic modulus $E$ & Weight of books on table \\
Elastic solid & Bulk modulus $K$ & Density of embedded inclusions \\
Planck field & $\rhostiff = c^2/(8\pi G)$ & $\rho_{\mathrm{matter}} \sim \rhocrit$ \\
\bottomrule
\end{tabular}
\end{table}

The stiffness density $\rhostiff$ is a \textbf{constitutive property} of the Planck field:
\begin{enumerate}[label=(\roman*)]
    \item It is an intensive quantity, independent of the volume under consideration.
    \item It is determined by fundamental constants ($c$, $G$) characterizing the field's intrinsic structure.
    \item It is unaffected by local matter content---$G$ is constant whether in intergalactic voids or near neutron stars (verified to $|\dot{G}/G| < 10^{-12}$ yr$^{-1}$).
\end{enumerate}

The cosmological critical density $\rhocrit = 3H_0^2/(8\pi G)$ is a \textbf{content property}---it describes the total matter, radiation, and dark energy distributed within the Planck field, set by the expansion rate $H_0$.

Comparing $\rhostiff$ to $\rhocrit$ is therefore a \textbf{category error}, analogous to comparing the tension in a guitar string to the mass of dust sitting on it.

\subsubsection{The Resting Tension}
\label{sec:T0_value}

The \emph{resting tension} of a medium is its stress-energy state in equilibrium, before external perturbations. For the Planck field:
\begin{equation}
\label{eq:T0_def}
T_0 \equiv K = \rhostiff \times c^2 = \frac{c^4}{8\pi G}.
\end{equation}

Numerically:
\begin{equation}
\label{eq:T0_numerical}
T_0 = \frac{(2.998 \times 10^8)^4}{8\pi \times 6.674 \times 10^{-11}} = 4.83 \times 10^{42} \text{ Pa}.
\end{equation}

This enormous resting tension---above nuclear pressure ($\sim 10^{34}$ Pa) but far below Planck pressure ($\sim 10^{113}$ Pa)---is the physical quantity that determines the gravitational constant through:
\begin{equation}
\label{eq:G_from_T0}
\boxed{G = \frac{c^4}{8\pi T_0}.}
\end{equation}

\begin{table}[H]
\centering
\caption{Comparison of Resting Tension $T_0$ to Known Pressure Scales}
\label{tab:pressure_scales}
\begin{tabular}{@{}llll@{}}
\toprule
\textbf{Pressure Scale} & \textbf{Value (Pa)} & \textbf{Ratio to $T_0$} & \textbf{Physical Context} \\
\midrule
Planck pressure $P_{\mathrm{Pl}}$ & $4.63 \times 10^{113}$ & $10^{71}$ & Quantum gravity limit \\
\textbf{Resting tension} $T_0$ & $4.83 \times 10^{42}$ & \textbf{1} & \textbf{Planck field stiffness} \\
Nuclear matter & $\sim 10^{34}$ & $10^{-8}$ & Neutron star cores \\
Center of the Sun & $\sim 10^{16}$ & $10^{-26}$ & Stellar interior \\
Dark energy $|P_\Lambda|$ & $\sim 10^{-10}$ & $10^{-52}$ & Cosmological scale \\
\bottomrule
\end{tabular}
\end{table}

\subsubsection{Why Gravity is Weak: The Fractional Displacement}
\label{sec:weak_gravity}

In the author's displacement mechanism \cite{SkavbergMechG2025,SkavbergMechGR2025}, gravity arises from matter \emph{displacing} the Planck field. The undisplaced field maintains resting tension $T_0$. When matter with energy density $\rho_{\mathrm{matter}} c^2$ is introduced, it perturbs the field:
\begin{equation}
\label{eq:perturbation}
\frac{\delta T}{T_0} \sim \frac{\rho_{\mathrm{matter}}}{\rhostiff}.
\end{equation}

For cosmological matter densities ($\rho_{\mathrm{matter}} \sim \rhocrit$):
\begin{equation}
\label{eq:weak_gravity_ratio}
\boxed{\frac{\delta T}{T_0} \sim \frac{\rhocrit}{\rhostiff} \sim 10^{-52}.}
\end{equation}

\textbf{Gravity is weak because matter barely perturbs the enormous resting tension of the Planck field.} The gravitational field arises as the restoring force of the field against this infinitesimal perturbation. The $10^{52}$ ratio quantifies this: matter creates a perturbation of only one part in $10^{52}$ of the background tension.

This interpretation explains the hierarchy problem of gravity: it is not that gravity is anomalously weak, but that the vacuum substrate is anomalously stiff.

\subsubsection{Resolution of the Density Discrepancy}
\label{sec:discrepancy_resolution}

The ratio of stiffness to critical density has a remarkably simple form:
\begin{equation}
\label{eq:density_ratio_main}
\frac{\rhostiff}{\rhocrit} = \frac{c^2/(8\pi G)}{3H_0^2/(8\pi G)} = \frac{c^2}{3H_0^2} = \frac{1}{3}\left(\frac{c}{H_0}\right)^2 = \frac{\RH^2}{3},
\end{equation}
where $\RH = c/H_0 \approx 1.37 \times 10^{26}$ m is the Hubble radius.

\begin{proposition}[Resolution of the Density Discrepancy]
The $\sim 10^{52}$ ratio $\rhostiff/\rhocrit = \RH^2/3$ is \textbf{not a fine-tuning problem}. It is a purely geometric factor encoding the macroscopic size of the observable universe. The ``discrepancy'' arises from comparing two fundamentally different quantities:
\begin{itemize}
    \item $\rhostiff$: how stiff the field is (constitutive)
    \item $\rhocrit$: how much matter is in the field (content)
\end{itemize}
A larger universe (larger $\RH$) naturally has a larger ratio.
\end{proposition}

\subsubsection{Implications for the Cosmological Constant Problem}
\label{sec:CC_implications}

The standard cosmological constant problem is the $\sim 10^{123}$ discrepancy between Planck density and observed dark energy density. The constitutive interpretation introduces an intermediate scale that factorizes this ratio:
\begin{equation}
\label{eq:CC_factorization}
\frac{\rhoPl}{\rhocrit} = \underbrace{\frac{\rhoPl}{\rhostiff}}_{\sim 10^{71}} \times \underbrace{\frac{\rhostiff}{\rhocrit}}_{\sim 10^{52}} = 10^{71} \times 10^{52} = 10^{123}.
\end{equation}

\textbf{Key insight:} The second factor ($10^{52}$) is \emph{not} a fine-tuning puzzle---it equals $\RH^2/3$ and reflects the geometric size of the universe. The first factor ($10^{71}$) represents the \emph{true} puzzle: why is the vacuum stiffness $10^{71}$ times below the Planck scale?

This interpretation reduces the cosmological constant problem from explaining $10^{123}$ to explaining $10^{71}$---a reduction of 52 orders of magnitude in the unexplained hierarchy.

\subsubsection{Constancy of $G$ from Constitutive Tension}
\label{sec:G_constancy}

Since $T_0$ is a constitutive property---determined by the Planck field's internal structure, not by external matter---$G$ is automatically constant:
\begin{equation}
G = \frac{c^4}{8\pi T_0} = \text{const}.
\end{equation}

This satisfies all observational bounds:
\begin{itemize}
    \item BBN: $|\Delta G/G| < 5\%$ over $\sim 10^{10}$ years.
    \item Lunar laser ranging: $|\dot{G}/G| < 10^{-13}$ yr$^{-1}$.
    \item Binary pulsars: $|\dot{G}/G| < 10^{-12}$ yr$^{-1}$.
\end{itemize}

The constancy of $G$ is not an additional assumption---it follows automatically from the constitutive nature of $T_0$.

\subsubsection{Falsifiable Prediction}
\label{sec:falsifiable_tension}

The constitutive interpretation provides a falsifiable relation:
\begin{equation}
\boxed{T_0 = \frac{c^4}{8\pi G} \quad \Leftrightarrow \quad G = \frac{c^4}{8\pi T_0}.}
\end{equation}

Any measurement of $G$ constrains $T_0$. If the resting tension could be independently measured (e.g., through quantum vacuum effects), the relation would be directly testable. Current precision on $G$ ($\sim 10^{-5}$ relative uncertainty) constrains $T_0$ to the same precision.

\subsubsection{Summary}
\label{sec:tension_summary}

The stiffness density $\rhostiff = c^2/(8\pi G)$ is a constitutive property of the Planck field, representing the resting tension:
\begin{equation}
\boxed{T_0 = \rhostiff c^2 = \frac{c^4}{8\pi G} = 4.83 \times 10^{42} \text{ Pa}.}
\end{equation}

\textbf{Key results:}
\begin{enumerate}
    \item $\rhostiff$ characterizes the \emph{medium itself}, not its matter content.
    \item The ``density discrepancy'' $\rhostiff/\rhocrit \sim 10^{52}$ is a \textbf{category error}: these are different physical quantities.
    \item Gravity is weak because $\delta T/T_0 \sim \rho_{\mathrm{matter}}/\rhostiff \sim 10^{-52}$---matter barely perturbs the field.
    \item The ratio $\rhostiff/\rhocrit = \RH^2/3$ is \textbf{geometric}, not a fine-tuning.
    \item Constant $T_0$ implies constant $G$, satisfying all observational bounds.
    \item The cosmological constant problem is reduced from $10^{123}$ to $10^{71}$.
\end{enumerate}

% ============================================================================
% END OF DRAFT-READY SECTION
% ============================================================================
